\documentclass{article}

% if you need to pass options to natbib, use, e.g.:
%     \PassOptionsToPackage{numbers, compress}{natbib}
% before loading neurips_2026

% The authors should use one of these tracks.
% Before accepting by the NeurIPS conference, select one of the options below.
% 0. "default" for submission
\usepackage[main]{neurips_2026}
% the "default" option is equal to the "main" option, which is used for the Main Track with double-blind reviewing.
% 1. "main" option is used for the Main Track
%  \usepackage[main]{neurips_2026}
% 2. "position" option is used for the Position Paper Track
%  \usepackage[position]{neurips_2026}
% 3. "eandd" option is used for the Evaluations & Datasets Track
 % \usepackage[eandd]{neurips_2026}
 % if you need to opt-in for a single-blind submission in the E&D track:
 %\usepackage[eandd, nonanonymous]{neurips_2026}
% 4. "creativeai" option is used for the Creative AI Track
%  \usepackage[creativeai]{neurips_2026}
% 5. "sglblindworkshop" option is used for the Workshop with single-blind reviewing
 % \usepackage[sglblindworkshop]{neurips_2026}
% 6. "dblblindworkshop" option is used for the Workshop with double-blind reviewing
%  \usepackage[dblblindworkshop]{neurips_2026}

% After being accepted, the authors should add "final" behind the track to compile a camera-ready version.
% 1. Main Track
 % \usepackage[main, final]{neurips_2026}
% 2. Position Paper Track
%  \usepackage[position, final]{neurips_2026}
% 3. Evaluations & Datasets Track
 % \usepackage[eandd, final]{neurips_2026}
% 4. Creative AI Track
%  \usepackage[creativeai, final]{neurips_2026}
% 5. Workshop with single-blind reviewing
%  \usepackage[sglblindworkshop, final]{neurips_2026}
% 6. Workshop with double-blind reviewing
%  \usepackage[dblblindworkshop, final]{neurips_2026}
% Note. For the workshop paper template, both \title{} and \workshoptitle{} are required, with the former indicating the paper title shown in the title and the latter indicating the workshop title displayed in the footnote.
% For workshops (5., 6.), the authors should add the name of the workshop, "\workshoptitle" command is used to set the workshop title.
% \workshoptitle{WORKSHOP TITLE}

% "preprint" option is used for arXiv or other preprint submissions
 % \usepackage[preprint]{neurips_2026}

% to avoid loading the natbib package, add option nonatbib:
%    \usepackage[nonatbib]{neurips_2026}

\usepackage[utf8]{inputenc} % allow utf-8 input
\usepackage[T1]{fontenc}    % use 8-bit T1 fonts
\usepackage{hyperref}       % hyperlinks
\usepackage{url}            % simple URL typesetting
\usepackage{booktabs}       % professional-quality tables
\usepackage{amsfonts}       % blackboard math symbols
\usepackage{nicefrac}       % compact symbols for 1/2, etc.
\usepackage{microtype}      % microtypography
\usepackage{xcolor}         % colors
\usepackage{amsmath,amsthm,stmaryrd}
\usepackage{algorithm}
\usepackage{algpseudocode}
\usepackage{listings}
\usepackage{tikz}
\usetikzlibrary{positioning}

\lstdefinelanguage{rust}{
  keywords={fn, let, mut, while, if, else, return, true, false, in, break, continue,
    Arc, Mutex, Condvar},
  sensitive=true,
  morecomment=[l]{//},
  morecomment=[s]{/*}{*/},
  morestring=[b]",
  morestring=[b]',
}
\usepackage{enumitem}
\usepackage{subcaption}
\usepackage{multirow}
\usepackage{graphicx}
\usepackage{pifont}
\newcommand{\cmark}{\ding{51}}
\newcommand{\xmark}{\ding{55}}
\setlength{\headheight}{15.5pt}
\emergencystretch=2em
\hbadness=4000
\vbadness=10000
\raggedbottom

% ── Listing style ──────────────────────────────────────────
\lstset{
  basicstyle=\ttfamily\footnotesize,
  keywordstyle=\bfseries,
  commentstyle=\color{gray},
  breaklines=true,
  frame=single,
  numbers=left,
  numberstyle=\tiny\color{gray},
  xleftmargin=1.6em,
  framexleftmargin=1.4em,
  columns=flexible,
  captionpos=b,
  aboveskip=6pt,
  belowskip=4pt
}

\newtheorem{definition}{Definition}
\newtheorem{theorem}{Theorem}
\newtheorem{proposition}{Proposition}
\newtheorem{lemma}{Lemma}
\newtheorem{assumption}{Assumption}
\newtheorem{remark}{Remark}

% Note. For the workshop paper template, both \title{} and \workshoptitle{} are required, with the former indicating the paper title shown in the title and the latter indicating the workshop title displayed in the footnote. 
\title{CIR+CVN: Bridging LLM Semantic Understanding and Petri-Net Verification for Concurrent Programs}

\author{%
  Kaiwen Zhang \\
  Department of Computer Science\\
  Tongji University\\
  Shanghai, China \\
  \texttt{zhangkw@tongji.edu.cn} \\
  \AND
  Guanjun Liu \\
  Department of Computer Science\\
  Tongji University\\
  Shanghai, China \\
  \texttt{liuguanjun@tongji.edu.cn} \\
}


\begin{document}
\maketitle
\begin{abstract}
Recovering concurrency structure directly from source code is difficult because shared-resource identity and protection relations are often obscured by aliasing, ownership, and API-specific idioms. We therefore study a specification-driven, model-first verification architecture for LLM-assisted concurrent program construction. Instead of verifying arbitrary source code, a large language model first synthesizes a verification-oriented concurrency artifact from a natural-language requirement or system specification. The first formalism, the Concurrency Intermediate Representation (\textsc{Cir}), is a statement-level, alias-free model in which shared resources are globally named, protection relations are explicit, and each statement carries a stable identifier. The second formalism, the Concurrency Verification Net (\textsc{Cvn}), is a weighted place/transition Petri net with a finite global store and three-valued guards for data-dependent branching. A validated \textsc{Cir} artifact is translated mechanically to \textsc{Cvn}, explored exhaustively, and any counterexample is mapped back to statement identifiers to guide targeted repair. To reduce the risk of bug-free but behavior-dropping repairs, acceptance additionally applies a lightweight goal-reachability check over designated critical outcomes. We formalize both representations, prove translation-correspondence results for deadlock and signal-loss analysis, define a two-layer checking architecture with 61 static rules and 5 analysis predicates, and evaluate the pipeline on 9 representative bounded-concurrency patterns. The results show that the method supports iterative bug detection and repair on \textsc{Cir} artifacts and that goal reachability helps filter semantically incomplete repairs. The trust boundary of the present work is the generated \textsc{Cir} artifact rather than arbitrary source code.
\end{abstract}

\section{Introduction}
\label{sec:introduction}

Concurrency bugs such as deadlocks, lost notifications, and blocking protocol errors remain difficult to detect and repair because they arise from interactions among threads rather than from a single local control path. Dynamic tools can expose concrete failing schedules but do not rule out unexplored interleavings~\cite{Godefroid1997VeriSoft,Flanagan2005DPOR,Musuvathi2008Heisenbugs,Burckhardt2010RandomizedScheduler}. Static techniques based on lock-order analysis, type systems, and ownership disciplines can provide broader guarantees~\cite{Engler2003RacerX,Flanagan2000TypeBasedRace,Flanagan2003Atomicity,Boyapati2002Ownership,Pratikakis2006Locksmith,Naik2006StaticRaceJava}, but they must first determine which operations refer to the same shared resource and which synchronization primitives protect which shared state. In practice this requires tracing values through heap allocation, aliasing, closure capture, and interprocedural flow. This reconstruction step is expensive, imprecise, and especially painful in automated repair workflows where diagnostics must be both precise and actionable.

Rather than verifying arbitrary source code, we ask whether concurrency verification can be moved earlier in the pipeline, to the stage where an LLM constructs a program from a natural-language specification. Recent work has shown that LLMs can find bugs, infer invariants, generate tests, and drive automated program repair~\cite{Chen2021CodeLLMs,Fan2023RepairFromLLMs,Xia2023APRLLMs,Bouzenia2025RepairAgent,Li2025ContextAwarePrompting}. However, LLMs alone cannot provide exhaustive reasoning about the combinatorial space of thread interleavings. Our work uses LLMs in a narrower way that LLM generates an alias-free concurrency model and later consumes structured counterexamples to repair it, while exhaustive bug discovery remains the responsibility of a formal verification engine. During construction the model can be asked to state the intended shared resources and synchronization structure explicitly, turning the problem from alias recovery into intent recovery.

This leads to our first formalism, the \emph{Concurrency Intermediate Representation}(\textsc{Cir}). \textsc{Cir} is a statement-level, alias-free concurrency model in which every shared resource has a unique global name, lock-to-variable protection relations are recorded explicitly, and every statement carries a stable identifier~($\mathsf{sid}$) for diagnostics and repair. A deterministic checker validates each generated artifact against structural and semantic rules before the more expensive verification phase.

The second formalism, the \emph{Concurrency Verification Net}(\textsc{Cvn}), provides the exhaustive interleaving analysis. Petri nets are a natural foundation for concurrency verification because causality, resource contention, and parallelism are represented directly in the net structure~\cite{Murata1989PetriNets,Esparza2008Unfoldings}. Standard P/T nets lack data-dependent branching. Classical colored Petri nets support full token-carried data but at considerable cost~\cite{Jensen2007CPNTools}. \textsc{Cvn} occupies a middle ground that it is a weighted P/T net augmented with a finite global variable store and three-valued guards. A guard that evaluates to \emph{unknown} causes both successor branches to be explored, preserving soundness through conservative over-approximation. Unlike prior Petri-net-based approaches that require manual modeling or language-specific extraction pipelines~\cite{Ball2002SLAM,Henzinger2003BLAST,Beyer2011CPAchecker}, \textsc{Cvn} is constructed mechanically from an LLM-generated \textsc{Cir} in which resource identity is already explicit. No alias reconstruction is needed at the net level.

A validated \textsc{Cir} is translated to \textsc{Cvn} and explored exhaustively. Any counterexample is mapped back to $\mathsf{sid}$ values and assembled into a structured diagnostic.  The LLM then uses this diagnostic to revise the affected synchronization logic, forming a generate-verify-repair loop~\cite{Clarke2000CEGAR,Bouzenia2025RepairAgent}. Acceptance requires both bug-freedom \emph{and} reachability of user's business goals. A repair that achieves safety by silently deleting required behavior is rejected. The trust boundary is the generated \textsc{Cir} artifact: we verify the model the LLM produces, not source code. Source-level extraction and conformance checking~\cite{Necula2000TranslationValidation, zhang2025petrinetsbasedmethodsautomatically} are complementary to our method. Extending the trust boundary to generated source code is an important direction for future work.

\medskip\noindent
The contributions of this paper are:
\begin{enumerate}[leftmargin=*,itemsep=2pt]

\item \textsc{Cir}, a verification-oriented concurrency intermediate representation with global resource naming, explicit protection mappings, and stable statement identifiers. It is embedded in a model-first generate-verify-repair paradigm with a lightweight goal-reachability filter that reduces the risk of semantically incomplete repairs.

\item \textsc{Cvn}, a weighted P/T Petri net with a finite global store and three-valued guards, together with a mechanical translation from \textsc{Cir}. We prove correspondence results showing that deadlock and signal-loss verdicts on the net faithfully reflect properties of the \textsc{Cir} artifact.

\item An experimental evaluation on 9~bounded-concurrency patterns and 5~LLMs demonstrating iterative bug detection and repair. The results show that goal-reachability checks filter semantically incomplete repairs that would otherwise pass safety verification.

\end{enumerate}

\section{Overview}
\label{sec:overview}

We introduce the method through a concrete example, then present the full pipeline.

\subsection{A signal-loss example}
\label{subsec:running-example}

Figure~\ref{fig:example-a} shows a Rust program with a signal-loss bug. The worker calls \texttt{wait} without a surrounding loop. The notifier calls \texttt{notify\_one} before setting the flag. If the notifier acquires the lock first, the signal fires with no waiter present and is lost. The worker then enters \texttt{wait} and blocks indefinitely.

\begin{figure}[t]
\begin{lstlisting}[language=rust]
fn worker(shared: Arc<(Mutex<bool>, Condvar)>) {
    let (m, cv) = &*shared;
    let g = m.lock().unwrap();
    let _g = cv.wait(g).unwrap(); // BUG: bare wait
}
fn notifier(shared: Arc<(Mutex<bool>, Condvar)>) {
    let (m, cv) = &*shared;
    let mut g = m.lock().unwrap();
    cv.notify_one();   // notify before write
    *g = true;
}
\end{lstlisting}
\caption{Signal-loss bug: bare wait without a loop guard.}
\label{fig:example-a}
\end{figure}

Figure~\ref{fig:cir-a} shows the corresponding \textsc{Cir} artifact. It is not extracted from source code but produced by an LLM from a natural-language specification. Alias analysis is unnecessary: \texttt{m0}, \texttt{cv0}, and \texttt{ready} are unique global names. The declaration \texttt{protection:\ ready:\ [m0]} makes the guarding relation explicit. Every statement carries a stable $\mathsf{sid}$ that anchors diagnostics across repair iterations.

\begin{figure}[t]
\begin{lstlisting}[basicstyle=\small\ttfamily, frame=single]
resources:
  m0:    { kind: Mutex }
  cv0:   { kind: Condvar, paired_with: m0 }
  ready: { kind: Var, type: Bool, init: false }
protection:  ready: [m0]
threads:
  worker:
    body:
    - { sid: w1, op: lock(m0),        next: w2 }
    - { sid: w2, op: wait(cv0, m0),   next: w3 }
    - { sid: w3, op: unlock(m0) }
  notifier:
    body:
    - { sid: n1, op: lock(m0),            next: n2 }
    - { sid: n2, op: notify(cv0),         next: n3 }
    - { sid: n3, op: write(ready, true),  next: n4 }
    - { sid: n4, op: unlock(m0) }
\end{lstlisting}
\caption{\textsc{Cir} for the signal-loss example.}
\label{fig:cir-a}
\end{figure}

The \textsc{Cir} artifact is translated mechanically into a \textsc{Cvn} and explored exhaustively. The verifier detects the signal-loss bug and emits the structured repair report shown in Figure~\ref{fig:repair-report}. The report identifies the bug kind, anchors the root cause to $\mathsf{sid}$ values \texttt{n2} and \texttt{w2}, and provides a concrete repair suggestion expressed in \textsc{Cir} constructs rather than Petri-net internals. The goal check additionally verifies that the worker can still reach its completion state after repair.  A fix that achieves safety by deleting behavior (e.g., removing all \texttt{notify} calls) would be rejected as behaviorally incomplete.

\begin{figure}[t]
\begin{lstlisting}[basicstyle=\small\ttfamily, frame=single]
bug_kind:  signal_loss
blame:     [n2, w2]
end_state:
  V[cv0_worker]: waiting
  V[ready]:      true
  held:          {}
  waiting:       [worker]
trace:     [n1, n2(lost), n3, n4, w1, w2(blocked)]
diagnosis:
  notify at sid=n2 fires when no thread is
  waiting on cv0; worker enters bare wait at
  sid=w2 with no loop guard and no future
  notification.
repair_suggestion:
  Insert read(ready) with branch before w2;
  add back-edge from post-wakeup to the condition
  check, forming a while(!ready) loop around
  wait(cv0, m0).
goal_check:
  critical_outcome: "worker exits synchronization"
  status: UNREACHABLE in current artifact
\end{lstlisting}
\caption{Structured repair report for the signal-loss example.}
\label{fig:repair-report}
\end{figure}

% ─────────────────────────────────────────────────────
\subsection{Pipeline and repair hierarchy}
\label{subsec:pipeline}

The architecture organizes the workflow into five stages, two verification layers, and three repair tiers.

% \begin{figure}[t]
% \centering
% \begin{tikzpicture}[
%   box/.style={rectangle, draw, rounded corners=2pt, thick,
%     text width=3.4cm, minimum height=0.7cm, align=center,
%     font=\footnotesize},
%   acc/.style={box, fill=black!8},
%   node distance=0.8cm
% ]

% % ── Nodes ──
% \node[box]               (s1) {Stage 1: \textsc{Cir} Generation};
% \node[box, below=of s1]  (s2) {Stage 2: Static Checking};
% \node[box, below=of s2]  (s3) {Stage 3: \textsc{Cir}$\;\to\;$\textsc{Cvn}};
% \node[box, below=of s3]  (s4) {Stage 4: State-Space Analysis};
% \node[box, below=of s4]  (s5) {Stage 5: Goal Reachability};
% \node[acc, below=of s5]  (ok) {Accepted};

% % ── Main flow ──
% \draw[->] (s1.south) -- (s2.north);
% \draw[->] (s2.south) -- (s3.north)
%   node[midway, right, font=\scriptsize] {pass};
% \draw[->] (s3.south) -- (s4.north);
% \draw[->] (s4.south) -- (s5.north)
%   node[midway, right, font=\scriptsize] {no bug};
% \draw[->] (s5.south) -- (ok.north)
%   node[midway, right, font=\scriptsize] {reachable};

% % ── Tier 1: auto-fix self-loop on Stage 2 ──
% \draw[->, dashed] ([yshift=2pt]s2.east)
%   .. controls +(0.7,0.35) and +(0.7,-0.35) ..
%   ([yshift=-2pt]s2.east);
% \node[font=\scriptsize] at ([xshift=10mm]s2.east) {Tier 1};

% % ── Tier 2: Stage 2 → Stage 1 (left) ──
% \draw[->, dashed] (s2.west) -- ++(-6mm,0) |- (s1.west);
% \node[font=\scriptsize, anchor=east]
%   at ([xshift=-6mm, yshift=4.5mm]s2.west) {Tier 2};

% % ── Tier 3: Stage 4 → Stage 1 (right, inner) ──
% \draw[->, dashed] (s4.east) -- ++(11mm,0)
%   |- ([yshift=-1mm]s1.east);
% \node[font=\scriptsize, anchor=west]
%   at ([xshift=11mm, yshift=5mm]s4.east) {Tier 3};

% % ── Goal fail: Stage 5 → Stage 1 (right, outer) ──
% \draw[->, dashed] (s5.east) -- ++(16mm,0)
%   |- ([yshift=1mm]s1.east);
% \node[font=\scriptsize, anchor=west]
%   at ([xshift=16mm, yshift=5mm]s5.east) {goal fail};

% \end{tikzpicture}
% \caption{The generate-verify-repair pipeline.}
% \label{fig:pipeline}
% \end{figure}

\textbf{Stage~1} generates a \textsc{Cir} artifact from a natural-language specification using an LLM.
\textbf{Stage~2} validates the artifact against 61~well-formedness rules organized in 9~categories, catching malformed or semantically inconsistent artifacts before the more expensive verification phase.
\textbf{Stage~3} translates the validated \textsc{Cir} into a \textsc{Cvn} through a stateless, deterministic translator. 
\textbf{Stage~4} explores the \textsc{Cvn} through exhaustive reachability search and reports two primary properties: deadlock and behavioral dead transition.  Signal loss and channel block are reported only as sub-classifications of an already-detected deadlock. This stage runs entirely offline and does not consume LLM budget.
\textbf{Stage~5} checks whether each user-declared business goal remains reachable.  An artifact is accepted only if it is both violation-free and goal-reachable.

Three repair tiers determine the response when a problem is found. \emph{Tier~1} applies deterministic rewrites for mechanical defects such as a missing \texttt{unlock} or a duplicate $\mathsf{sid}$, without involving the LLM. \emph{Tier~2} sends structural problems such as undefined resources or broken spawn/join topology to the LLM for localized regeneration. \emph{Tier~3} constructs a structured diagnostic from a concurrency violation or an unmet goal, including blame anchors, end-state summary, and repair suggestion, and asks the LLM to revise the affected synchronization logic. The loop iterates until the artifact passes all stages or a configurable budget (default: 5~rounds) is exhausted.


\section{Method}
\label{sec:method}

This section defines the two core formalisms (\textsc{Cir} and \textsc{Cvn}), describes the translation between them, and presents the bug-detection and repair machinery.

\subsection{Concurrency intermediate representation}
\label{subsec:cir-def}

A \textsc{Cir} artifact is a statement-level concurrency model designed to be generated by an LLM and consumed by a formal analysis engine.

\begin{definition}[\textsc{Cir} artifact]
\label{def:cir}
A \textsc{Cir} artifact is a 6-tuple $\mathcal{I}=(R,\;G,\;F,\;\Sigma,\;e,\;\Gamma)$, where $R$ is a finite resource table mapping names to kinds and
configurations, $G:\mathit{VarName}\rightharpoonup\mathcal{P}(\mathit{LockName})$ is a protection mapping, $F$ is a finite set of function definitions, $\Sigma$ is a partial function-summary map for external calls, $e\in F$ is the designated entry function, and $\Gamma$ is a finite set of business goals.
\end{definition}

Every resource in $R$ has a unique global name. Pointers, heap indirection, and aliasing are strictly forbidden. Functions reference resources directly by name.  Each function body is a sequence of \emph{statements} $s=(\mathit{sid},\,\mathit{op},\,\mathit{transfer})$, where $\mathit{sid}$ is a globally unique identifier, $\mathit{op}$ is a concurrency or data operation (lock, wait, notify, send, recv, read, write, spawn, join, etc.), and $\mathit{transfer}$ specifies the control-flow successor. The complete operation catalogue and statement grammar are given in Appendix~\ref{app:cir-details}.

The protection mapping $G$ records which locks guard which variables. It is used exclusively for static checking and does not produce \textsc{Cvn} structure.  The summary map $\Sigma$ provides effect summaries for external functions whose bodies are not modeled. The static checker validates each artifact against 61~well-formedness rules organized in 9~categories before the verification phase.

A \emph{business goal} $\gamma=(C_\gamma,\,V_\gamma)$ specifies what the program is supposed to achieve: $C_\gamma$ contains completion and availability requirements (e.g., function $f$ must have returned, resource $r$ must be free), and $V_\gamma$ specifies required variable values. Goals are expressed entirely in \textsc{Cir} vocabulary. A repair is accepted only if no concurrency bug is detected \emph{and} every goal is reachable.

% ═══════════════════════════════════════════════════════
\subsection{Concurrency verification net}
\label{subsec:cvn-def}

\begin{definition}[\textsc{Cvn}]
\label{def:cvn}
A Concurrency Verification Net is an 8-tuple $\mathcal{N}=(P,\;T,\;W_{\!\mathit{in}},\;W_{\!\mathit{out}},\;G,\;U,\;\mathit{Var},\;\mu)$, where $P=P_c\uplus P_r\uplus P_a$ is a finite set of places partitioned into control places, resource places, and auxiliary places, $T$ is a finite set of transitions, $W_{\!\mathit{in}}$ and $W_{\!\mathit{out}}$ are input and output weight functions, $G:T\to\mathit{BoolExpr}$ assigns a guard to each transition, $U:T\to(\mathit{Var}\rightharpoonup\mathit{Expr})$ assigns variable updates, $\mathit{Var}$ is a finite set of variable names, and $\mu:T\to\mathit{SID}$ is an anchor mapping that traces every transition back to its originating \textsc{Cir} statement.
\end{definition}

A \textsc{Cvn} state is a pair $S=(M,V)$ of a marking $M:P\to\mathbb{N}_0$ and a valuation $V:\mathit{Var}\to\mathit{Val}$. The value domain extends concrete
values with an absorbing element $\top$ (unknown): any comparison involving $\top$ evaluates to the truth value $\mathsf{unknown}$. A transition $t$ is enabled at $S$ iff every input place has enough tokens and its guard does not evaluate to $\mathsf{false}$. When a guard evaluates to $\mathsf{unknown}$, the transition is treated as enabled.  For a branch, both successors are explored simultaneously. This over-approximation ensures that no reachable state is missed. The standard Petri-net firing rule applies to tokens. Variable updates are evaluated against the pre-firing valuation. The complete firing semantics, place taxonomy, and initial token assignments are given in Appendix~\ref{app:cvn-details}.

\subsection{Translation from CIR to CVN}
\label{subsec:translation}
Let $\textsc{Cir}_v$ denote the subset of \textsc{Cir} artifacts that pass all static-checking rules. The translator is a pure, total function $\mathit{translate}:\textsc{Cir}_v\to\textsc{Cvn}$ executed in three phases. \textbf{Phase~1} scans the resource table and generates resource places with their initial token counts (e.g., Mutex $\mapsto 1$, Semaphore$(n)\mapsto n$). \textbf{Phase~2} translates each function body: every statement produces at least one anchored transition, preserving diagnostic traceability. \textbf{Phase~3} translates function summaries: each summarized call becomes a single transition that sets all written variables to $\top$. The translation also maps each business goal to a \textsc{Cvn} reachability query. A state $(M,V)$ \emph{satisfies} a goal $\gamma$ iff all completion and availability requirements hold in $M$ and all required variable values match in $V$. A goal is \emph{reachable} if at least one reachable state satisfies it. The full translation rules are given in Appendix~\ref{app:translation-rules}.

\subsection{Bug detection}
\label{subsec:bug-detection}

Given a translated \textsc{Cvn} $(\mathcal{N},S_0)$, the verifier performs exhaustive state-space exploration from $S_0$. Two primary properties are reported.

\textbf{Deadlock.}\;
A state $S$ is a deadlock iff no transition is enabled and at least one control token remains on a non-return place. Every such state is reported with its full firing sequence from $S_0$. Two refinements can further classify a deadlock witness. A \emph{signal-loss} refinement applies when blocked tokens reside on condvar wait-places or when the witness contains a notification that fired with no waiter present. A \emph{channel-block} refinement applies when a blocked
consumer requires tokens from an empty channel place. Both refinements are applied strictly after a deadlock has been confirmed, so they never enlarge the reported bug set.

\textbf{Behavioral dead transition.}\;
A transition $t$ is behaviorally dead iff it never fires on any edge of the reachability graph. Such a transition corresponds to a \textsc{Cir} statement that can never execute under any feasible interleaving.  This property is computed from the same reachability graph used for deadlock detection, so it requires no additional search budget. 

The verifier deliberately avoids scheduler-dependent properties such as livelock and starvation, which cannot be decided soundly on scheduler-agnostic Petri-net semantics.

\subsection{Goal reachability}
\label{subsec:goal-reach}

A \textsc{Cir} that passes the deadlock check may still fail to implement its declared behavior. We promote business-goal reachability to a first-class property checked with the same state-space engine. Each goal $\gamma\in\Gamma$ is compiled into a conjunctive predicate over \textsc{Cvn} states. A goal is
\emph{reachable} iff some reachable state simultaneously satisfies all its predicates. The verifier returns the set $\mathcal{U}\subseteq\Gamma$ of unmet goals. Because $\top$ never equals a concrete value, unresolved branches cannot spuriously satisfy a goal. The formal soundness statement is given in Theorem~\ref{thm:goal-soundness}.

\subsection{Goal-aware repair loop}
\label{subsec:repair-loop}

The repair loop implements the three-tier hierarchy. Algorithm~\ref{alg:repair-goal} gives the process.

\begin{algorithm}[t]
\caption{Goal-Aware Generate-Verify-Repair Loop}
\label{alg:repair-goal}
\begin{algorithmic}[1]
\Require Specification $\Psi$, iteration budget $K$
\Ensure  Verified \textsc{Cir} $\mathcal{I}$ or Failure
\State $\mathcal{I} \gets \text{LLM.Generate}(\Psi)$
\For{$\mathit{round}=1$ \textbf{to} $K$}
    \State $\mathcal{E} \gets \text{StaticCheck}(\mathcal{I})$
    \If{$\mathcal{E} \neq \emptyset$}
        \State Apply Tier-1 or Tier-2 repair; \textbf{continue}
    \EndIf
    \State $\mathcal{N} \gets \text{Translate}(\mathcal{I})$
    \State $(\mathcal{R},\,\mathcal{B})
      \gets \text{Analyze}(\mathcal{N})$
    \If{$\mathcal{B} \neq \emptyset$}
        \State $\mathcal{D} \gets
          \text{BuildDiag}(\text{SelectBug}(\mathcal{B}),
            \mathcal{R},\mathcal{I})$
        \State $\mathcal{I} \gets
          \text{LLM.Repair}(\mathcal{I},\mathcal{D})$;
          \textbf{continue}
    \EndIf
    \State $\mathcal{U} \gets
      \text{CheckGoals}(\mathcal{N},\Gamma,\mathcal{R})$
    \If{$\mathcal{U}=\emptyset$}
        \Return $\mathcal{I}$
    \EndIf
    \State $\mathcal{I} \gets
      \text{LLM.Repair}\bigl(\mathcal{I},\,
        \text{BuildGoalDiag}(\mathcal{U})\bigr)$
\EndFor
\State \Return Failure
\end{algorithmic}
\end{algorithm}

$\text{Analyze}$ returns the reachable state set $\mathcal{R}$ and the bug set $\mathcal{B}$, which contains detected deadlocks (already refined to signal-loss or channel-block where applicable) and behavioral dead transitions. When multiple bugs are found, $\text{SelectBug}$ picks one by priority ($\text{deadlock}>\text{signal-loss}>\text{channel-block}> \text{dead-transition}$), preferring the shortest witness.

The key ordering invariant is that goal reachability is checked only after the bug set is empty. If a goal is unreachable because a deadlock blocks its producer, the deadlock is repaired first. This gives every round a single semantic target: either remove a bug or restore missing behavior.

When a bug is found, the verifier constructs a structured diagnostic containing the bug kind, a witness trace expressed in $\mathsf{sid}$ values, a state summary (thread positions, held resources, waiting relations), a bounded \textsc{Cir} slice around implicated statements, preservation constraints, and a repair hint.  The diagnostic provides both \emph{causality} (how the system reached the bug) and \emph{locality} (where it is stuck). The full diagnostic format and
prompt assembly are described in Appendix~\ref{app:repair-details}.

\section{Evaluation}
\label{sec:evaluation}

\subsection{Setup}
\label{subsec:setup}

\paragraph{Benchmark.}
Table~\ref{tab:test-matrix} lists 9~bounded-concurrency patterns used throughout the evaluation. Six patterns contain known defects: five are definite concurrency bugs and one (\textsc{partial\_deadlock}) is a behavioural regression in which two threads stall each other while a bystander keeps the system globally live. The bystander prevents a global deadlock, so the defect is visible only through the goal-reachability layer. Three patterns are bug-free baselines used to confirm the absence of false positives.

\begin{table*}[t]
\caption{Benchmark patterns.
$|\Gamma|$: number of business goals checked after repair.}
\label{tab:test-matrix}
\centering\footnotesize
\begin{tabular}{@{}clll c p{3.6cm}@{}}
\toprule
\textbf{\#} & \textbf{Pattern} & \textbf{Bug kind}
  & \textbf{Typical repair}
  & $|\Gamma|$
  & \textbf{Goal summary} \\
\midrule
1  & Two-mutex deadlock & Deadlock
   & Unify lock order
   & 1 & Both threads ret; mutexes released \\[2pt]
2  & Condvar signal loss
   & SignalLoss 
   & Write pred.\ before notify
   & 2 & Worker ret; \texttt{V[ready]=true} \\[2pt]
3  & Channel\,+\,mutex DL
   & Deadlock
   & Release lock before ch.\ op
   & 1 & Both threads ret; mutex released \\[2pt]
4  & Three-lock circular
   & Deadlock 
   & Global lock order
   & 1 & All three threads ret \\[2pt]
5  & Partial deadlock
   & GoalUnreach$^\dagger$ 
   & Unify lock order
   & 2 & Thread\,A ret; Thread\,B ret \\[2pt]
6  & Dual condvar cross
   & Deadlock 
   & Break circular wait
   & 2 & Thread\,A ret; Thread\,B ret \\[2pt]
\midrule
7  & Semaphore throttle 
   & (none) 
   & --- & --- & --- \\[2pt]
8  & CAS contention 
   & (none) 
   & --- & --- & --- \\[2pt]
9  & FnSummary prop.\ 
   & (none)
   & --- & --- & --- \\
\bottomrule
\end{tabular}
\end{table*}

We evaluate five LLMs: GPT-5, Claude~4.7~Opus, Gemini~3~Pro, Qwen3-Max, and DeepSeek-V3.  All models use temperature\,0 with a 4\,096-token output limit. For each pattern a short natural-language concurrency specification serves as the generation input. Buggy and fixed Rust programs are maintained as reference implementations but are \emph{not} input to \textsc{Cir} generation. Generation and repair budgets are each capped at 5~rounds.  All experiments run on
Apple~M4~Pro\,/\,24\,GB; times are wall-clock averages over 3~runs.

\subsection{Translation Correctness and State-Space Scale}
\label{subsec:translation-results}

We verify two aspects of the \textsc{Cir}-to-\textsc{Cvn} translation: structural invariants and bug-detection accuracy. Nine invariants are checked on every translated \textsc{Cvn} (e.g., every CIR statement produces $\geq 1$ transition, initial marking matches resource declarations, spawn/join arcs are well-formed). All 81 checks (9~invariants $\times$ 9~patterns) pass.

Table~\ref{tab:translation} reports the translation size and bug-detection result for each pattern. Each buggy pattern uses a fixed ground-truth \textsc{Cir}, making results model-independent. All 5~definite bugs are detected with the correct bug kind. Pattern~5 contains no definite bug: the bystander thread keeps the system live, so the deadlock search correctly reports no deadlock. Instead, the goal-reachability layer flags that neither contending thread can reach its return place.  The 3~baselines produce zero false positives.

\begin{table}[t]
\caption{Translation size and detection result on ground-truth \textsc{Cir}.}
\label{tab:translation}
\centering\footnotesize
\begin{tabular}{@{}clcccrl@{}}
\toprule
\textbf{\#} & \textbf{Pattern}
  & \textbf{Stmts} & $|P|$ & $|T|$
  & \textbf{States} & \textbf{Result} \\
\midrule
1  & Two-mutex DL       & 15 & 22 & 17 &  60  & Deadlock \\
2  & Condvar SL         & 15 & 24 & 21 &  47  & SignalLoss \\
3  & Channel+mutex      & 13 & 20 & 15 &  37  & ChannelBlock \\
4  & Three-lock circ.   & 22 & 32 & 25 & 396  & Deadlock \\
5  & Partial DL         & 22 & 33 & 25 &  81  & GoalUnreach \\
6  & Dual condvar       & 19 & 32 & 29 &  21  & SignalLoss \\
\midrule
7  & Semaphore (BL)     & 16 & 24 & 19 & 187  & Verified \\
8  & CAS (BL)           & 11 & 16 & 15 &  31  & Verified \\
9  & FnSummary (BL)     & 15 & 21 & 17 &  57  & Verified \\
\bottomrule
\end{tabular}
\end{table}

All patterns remain below 400~reachable states. The largest (Pattern~4, three threads with three mutexes) completes analysis in under 1\,ms.  The average ratio is 1.6~places and 1.3~transitions per \textsc{Cir} statement. Condvar patterns exhibit a higher transition ratio ($\sim$1.5) due to the wait/notify decomposition. Goal-reachability checking adds negligible overhead by reusing the already-computed reachable set. Moving to $k \geq 4$ threads or richer data domains will increase the state space exponentially; partial-order reduction and symmetry breaking are planned future work.

\subsection{Generation and Repair}
\label{subsec:gen-repair}

Table~\ref{tab:gen-repair} reports two numbers per cell: the rounds needed to produce a statically valid \textsc{Cir} (generation) and the subsequent rounds needed to reach an accepted artifact (repair). Acceptance requires static validity, bug-freedom in the \textsc{Cvn} search, \emph{and} reachability of all declared business goals.

\begin{table*}[t]
\caption{Generation rounds / repair rounds per model per pattern.}
\label{tab:gen-repair}
\centering\footnotesize
\begin{tabular}{@{}cl ccccc@{}}
\toprule
\textbf{\#} & \textbf{Pattern}
  & \textbf{GPT-5} & \textbf{Claude 4.7 Opus}
  & \textbf{Gemini 3 Pro}
  & \textbf{Qwen3-Max} & \textbf{DeepSeek-V3} \\
\midrule
1  & Two-mutex DL
   & 1\,/\,1 & 1\,/\,1 & 1\,/\,2
   & 1\,/\,1 & 2\,/\,1 \\
2  & Condvar SL
   & 1\,/\,5$^{+4}$ & 1\,/\,1 & 2\,/\,1
   & 1\,/\,\ding{55}$^{+5}$ & 3\,/\,\ding{55}$^{+5}$ \\
3  & Channel+mutex
   & 1\,/\,1 & 1\,/\,1 & 1\,/\,1
   & 2\,/\,1 & 2\,/\,1 \\
4  & Three-lock circ.
   & 2\,/\,1 & 1\,/\,1 & 2\,/\,1
   & 2\,/\,1 & 3\,/\,1 \\
5  & Partial DL
   & 2\,/\,5 & 1\,/\,1 & 2\,/\,2
   & 3\,/\,2$^{+1}$ & 4\,/\,\ding{55} \\
6  & Dual condvar
   & 2\,/\,\ding{55}$^{+2}$ & 2\,/\,\ding{55}$^{+4}$
   & 3\,/\,3$^{+2}$
   & 3\,/\,\ding{55}$^{+5}$ & \ding{55}\,/\,---  \\
\midrule
7  & Semaphore (BL)
   & 1\,/\,--- & 1\,/\,--- & 1\,/\,---
   & 1\,/\,--- & 1\,/\,--- \\
8  & CAS (BL)
   & 1\,/\,--- & 1\,/\,--- & 1\,/\,---
   & 2\,/\,--- & 2\,/\,--- \\
9  & FnSummary (BL)
   & 1\,/\,--- & 1\,/\,--- & 1\,/\,---
   & 2\,/\,--- & 2\,/\,--- \\
\midrule
\multicolumn{2}{@{}l}{Gen.\ success}
   & 9/9 & 9/9 & 9/9 & 9/9 & 8/9 \\
\multicolumn{2}{@{}l}{Repair accepted}
   & 5/6 & 5/6 & 6/6 & 4/6 & 3/6 \\
\multicolumn{2}{@{}l}{Regressions}
   & 6 & 4 & 2 & 11 & 7 \\
\bottomrule
\end{tabular}
\end{table*}

\paragraph{Generation.}
$\ding{55}=$ not achieved within 5~rounds. Superscript~${}^{+n}$: intermediate regressions caught by the back-end. Baselines (Patterns 7--9) require generation only that no bugs to repair. All models except DeepSeek-V3 produce a valid \textsc{Cir} for every pattern.  DeepSeek-V3 fails on Pattern~6, the most structurally complex pattern requiring two condvars, two mutexes, and cross-thread notification dependencies. Claude~4.7~Opus achieves the lowest average generation cost (1.1~rounds), with first-attempt success on 8 of 9~patterns. The most common generation errors are malformed condvar declarations, incorrect branch conditions, and omitted protection mappings. All are caught by the static checker and typically corrected within one feedback round.

\paragraph{Bug fidelity.}
Before entering repair, we verify that each model's first valid \textsc{Cir} actually contains the intended bug. GPT-5, Claude~4.7~Opus, and Gemini~3~Pro achieve 6/6 fidelity: every generated \textsc{Cir} exhibits exactly the expected bug kind. Two fidelity failures occur in the remaining models. Qwen3-Max on Pattern~5 generates both threads with the same lock order, eliminating the intended conflict. DeepSeek-V3 on Pattern~4 similarly unifies lock acquisition orders, breaking the circular dependency. These cases are harmless for end users (the \textsc{Cir} is already correct), but the repair loop is never exercised for that model-pattern pair. No case arises in which a model captures a bug of a \emph{different} kind than specified.

\paragraph{Repair.}
Gemini~3~Pro is the only model that repairs all 6~defective patterns, with 5/6 accepted on the first attempt. Claude~4.7~Opus achieves single-round repair on every accepted pattern (5/6) but fails on the dual-condvar cross-wait (Pattern~6). GPT-5 requires five rounds on Pattern~2 (signal loss) and Pattern~5 (goal-unreachability). Qwen3-Max and DeepSeek-V3 succeed on the three simple-deadlock patterns but struggle with condvar-heavy patterns and the goal-reachability pattern.

Pattern~6 is the hardest cell: only Gemini~3~Pro disentangles the two cross-notification dependencies without introducing fresh bugs. Common failure modes include moving a \textsf{wait} off its associated mutex and introducing a new signal-loss by emitting \textsf{notify} outside the protecting critical section.  These regressions are caught by the back-end and surfaced to the LLM, but five rounds are not enough for the other models to converge.

Across the 30 model-pattern repair tasks, 30~intermediate regressions occur (static errors: $\sim$13, translation violations: $\sim$14, goal-unreachability: $\sim$1, bug-kind drift: $\sim$2).  No regression survives to a final accepted \textsc{Cir}.

\subsection{Goal-Reachability Effectiveness}
\label{subsec:goal-results}

Goal-reachability checking serves two roles: it catches repair-induced regressions that silently drop functionality, and it detects partial deadlocks that leave the system globally live but permanently block the threads responsible for reaching declared goals.

On the ground-truth \textsc{Cir}s, the checker correctly reports all goals as met on fixed patterns and all goals as unmet on buggy patterns (4/4 verdicts correct on the two goal-bearing patterns, 5 and~9). On bug-free baselines no false alarms arise.

The more revealing test is the checker's role inside the repair loop. Pattern~5 is invisible to the deadlock search because a bystander thread keeps the system live. Without goal feedback the loop would receive no signal after the first bug-free candidate. With goal feedback, 4 of 5~models converge: Claude~4.7~Opus in 1~round, Gemini~3~Pro in 2, Qwen3-Max in 2, and GPT-5 in 5. DeepSeek-V3 fails to rearrange the control flows within budget. No accepted \textsc{Cir} in the entire evaluation has an unmet goal.

Table~\ref{tab:summary} consolidates the results. Acceptance requires static validity, bug-freedom, and goal reachability.

\begin{table}[t]
\caption{End-to-end summary across the 6 defective patterns.}
\label{tab:summary}
\centering\footnotesize
\begin{tabular}{@{}l ccccc@{}}
\toprule
& \textbf{GPT-5} & \textbf{Cl.~4.7}
  & \textbf{Gem.~3}
  & \textbf{Qwen3} & \textbf{DS-V3} \\
\midrule
Bug-free (P1--4,6)
             & 4/5 & 4/5 & 5/5 & 3/5 & 3/5 \\
Goal reach.\ (P5)
             & \ding{51} & \ding{51} & \ding{51}
             & \ding{51} & \ding{55} \\
Total repair rounds
             & 13   & 5   & 10   & 5   & 3   \\
Regressions
             & 6    & 4   & 2    & 11  & 7   \\
\midrule
\textbf{Accepted}
             & 5/6 & 5/6 & 6/6 & 4/6 & 3/6 \\
\bottomrule
\end{tabular}
\end{table}

Only Gemini~3~Pro is accepted on all 6~patterns. The hardest pattern (dual condvar cross-wait) requires coordinating wait/notify placements across two threads simultaneously. Pattern~5 (partial deadlock) is visible only through goal-reachability feedback. Once that feedback is available, 4 of 5~models converge.  Across the entire matrix the pipeline never accepts a \textsc{Cir} that contains a definite concurrency bug or whose declared business goals are unreachable. Every intermediate regression is caught by one of the three verification layers: static checking, \textsc{Cvn} analysis, or goal reachability.


\section{Conclusion}
\label{sec:conclusion}

This paper shifts concurrency verification from source-level alias recovery to a closed loop over \textsc{Cir}, an explicit, alias-free concurrency model constructed by LLMs from natural-language specifications. A deterministic translation to \textsc{Cvn} provides sound and complete detection of deadlock and signal loss under concrete valuation; livelock is flagged as an advisory structural warning with a contrapositive impossibility guarantee. Goal-reachability checking catches semantic regressions that escape all other verification layers. Across 9~patterns and 5~LLMs, no \textsc{Cir} containing a definite bug or violating a business goal is ever accepted.

% ═══════════════════════════════════════════════════════════════
% APPENDIX
% ═══════════════════════════════════════════════════════════════
\appendix

\section{CIR Details}
\label{app:cir-details}

\begin{definition}[Function and statement]
\label{def:statement}
A \emph{function definition} is a triple $(f,\,k,\,B)$ where $f$ is
the name,
$k\in\{\textsf{normal},\textsf{async},\textsf{closure}\}$ is the
kind, and $B=(s_1,\dots,s_n)$ is a body of statements.
A \emph{statement} is a triple
$s=(\mathit{sid},\;\mathit{op},\;\mathit{transfer})$, where
$\mathit{sid}\in\mathit{SID}$ is a unique identifier, $\mathit{op}$ is
an operation from Table~\ref{tab:cir-ops} or the distinguished
no-operation $\textsf{nop}$, and $\mathit{transfer}$ specifies the
control-flow successor:
\[
  \mathit{transfer} \in
  \bigl\{\,
    \textsf{next}(s'),\;
    \textsf{branch}(\beta,s_t,s_f),\;
    \textsf{switch}(x,\langle v_i\!\to\!s_i\rangle_{i},s_d),\;
    \textsf{return}
  \,\bigr\}
\]
where $\beta\in\mathit{BoolExpr}$ (Definition~\ref{def:expressions}).
The operation and transfer are orthogonal: $\mathit{op}$ specifies the
resource or concurrency action, and $\mathit{transfer}$ specifies the
successor.
\end{definition}

$\mathit{SID}$ denotes the universe of statement identifiers.  Every
function $(f,k,B)$ has a distinguished \emph{entry identifier}
$\mathit{sid}^{0}_{f}$ (the $\mathit{sid}$ of the first statement
in~$B$) and a distinguished \emph{return identifier}
$\mathit{ret}_{f}\notin\{\mathit{sid}\mid(\mathit{sid},\_,\_)\in B\}$.
When unambiguous we write $\mathit{sid}^{0}$ and $\mathit{ret}$
without the function subscript.  The sentinel $\mathit{sid}_{\bot}$ is
reserved for synthetic transitions generated from function summaries.

\begin{table}[h]
\caption{\textsc{Cir} operations.}
\label{tab:cir-ops}
\centering\footnotesize
\begin{tabular}{@{}llp{3.2cm}@{}}
\toprule
\textbf{Category} & \textbf{Operation} & \textbf{Target} \\
\midrule
\multirow{2}{*}{Lock}
  & \textsf{lock}, \textsf{drop}
  & Mutex, RwLock (write) \\
  & \textsf{read}, \textsf{drop}
  & RwLock (read) \\[3pt]
\multirow{3}{*}{Synchronization}
  & \textsf{wait}, \textsf{notify}, \textsf{notify\_all}
  & Condvar \\
  & \textsf{acquire}, \textsf{release}
  & Semaphore \\
  & \textsf{send}, \textsf{recv}
  & Channel \\[3pt]
\multirow{3}{*}{Data}
  & \textsf{read}, \textsf{write}
  & Var \\
  & \textsf{load}, \textsf{store}
  & Atomic \\
  & $\textsf{cas}(\mathit{expected},\mathit{new})$
  & Atomic \\[3pt]
\multirow{3}{*}{Control}
  & \textsf{spawn}, \textsf{join}
  & function (OS thread) \\
  & \textsf{spawn\_async}, \textsf{await}
  & function (async task) \\
  & \textsf{call}
  & function \\
\bottomrule
\end{tabular}
\end{table}

The summary map $\Sigma$ provides a summary for each function whose
body is not modeled in $F$.  A summary is a tuple
$\Sigma(f)=(\mathit{reads},\mathit{writes},\mathit{calls},
\mathit{has\_concurrency})$, where $\mathit{reads}$ and
$\mathit{writes}$ are sets of resource names, $\mathit{calls}$ is a
set of callee names, and $\mathit{has\_concurrency}$ is a boolean.
Only the $\mathit{writes}$ field affects the \textsc{Cvn} translation.
The remaining fields serve the static checker.

\begin{definition}[\textsc{Cir} business goal]
\label{def:cir-goal}
A \textsc{Cir} business goal is a pair $\gamma=(C_\gamma,\,V_\gamma)$,
where $C_\gamma$ is a set of \emph{completion} and \emph{availability}
requirements: $(f,\textsf{completed})$ requires function $f$ to have
returned, and $(r,\textsf{available})$ requires resource $r$ to be in
its initial (free) state.
$V_\gamma:\mathit{VarName}\rightharpoonup\mathit{Val}$ specifies
required variable values.  Goals are expressed entirely in \textsc{Cir}
vocabulary and do not reference \textsc{Cvn} places or markings.
\end{definition}

\begin{figure}[h]
\begin{lstlisting}[basicstyle=\small\ttfamily,frame=single]
goals:
  - id: G1
    desc: "Both threads complete, lock released,
           ready flag set"
    completion:
      - [worker, completed]
      - [notifier, completed]
    availability:
      - [m0, available]
    variables:
      ready: true
  - id: G2
    desc: "Worker reaches completion state"
    completion:
      - [worker, completed]
\end{lstlisting}
\caption{Business goals expressed in \textsc{Cir} vocabulary.}
\label{fig:goals-example}
\end{figure}

Business goals capture what the program is \emph{supposed to achieve},
not merely what it should avoid.  A buggy program may fail to reach
its goals because a deadlock or signal loss blocks the execution path.
Conversely, a repair that eliminates all concurrency bugs may still
fail to reach a goal if it inadvertently removes a critical operation.
The acceptance criterion for a verified \textsc{Cir} is therefore the
conjunction: (1)~no concurrency bug is detected, and (2)~every
business goal is reachable.  Goal reachability is checked only after
bug-freedom is established, because bugs themselves can block
reachability.  Goals can be specified in two ways: the user can state
them explicitly as part of the natural-language specification, or the
LLM can derive them automatically when generating the \textsc{Cir}
artifact.  In either case the goals become part of $\mathcal{I}$ and
travel through the entire pipeline.

% ═══════════════════════════════════════════════════════════════
\section{CVN Details}
\label{app:cvn-details}

\begin{definition}[Value domain]
\label{def:val}
The value domain is
$\mathit{Val}=\mathit{Concrete}(\tau)\cup\{\top\}$, where $\tau$
ranges over the base types
$\{\textsf{Bool},\textsf{Int},\textsf{Float},\textsf{String},
\textsf{Enum}\}$.  The element $\top$ (\emph{unknown}) is absorbing:
any arithmetic operation receiving $\top$ as an operand
produces $\top$, and any comparison receiving $\top$ produces the
truth value $\mathsf{unknown}$.  A literal write overwrites $\top$
unconditionally.
\end{definition}

\begin{definition}[Expressions]
\label{def:expressions}
\emph{Value expressions} and \emph{boolean expressions} over the value
domain are defined by:
\begin{align*}
  \mathit{Expr} &::=
    \mathit{Lit}(v) \mid
    \mathit{Ref}(x) \mid
    \mathit{BinOp}(\mathit{Expr},\oplus,\mathit{Expr}) \\[4pt]
  \mathit{BoolExpr} &::=
    \textsf{True} \mid \textsf{False} \mid
    \mathit{Cmp}(\mathit{Expr},\mathit{cmp},\mathit{Expr}) \\
    &\quad\mid\;
    \mathit{And}(\mathit{BoolExpr},\mathit{BoolExpr}) \mid
    \mathit{Or}(\mathit{BoolExpr},\mathit{BoolExpr}) \\
    &\quad\mid\;
    \mathit{Not}(\mathit{BoolExpr})
\end{align*}
Given a valuation $V$, $\mathit{eval}(e,V)$ evaluates a value
expression to an element of $\mathit{Val}$.  $\mathit{Ref}(x)$
evaluates to $V(x)$ and $\mathit{BinOp}$ follows the absorbing rule
for $\top$.
\end{definition}

\begin{definition}[Guard evaluation]
\label{def:guard-eval}
A \emph{guard} is a $\mathit{BoolExpr}$.  Given a valuation $V$, the
evaluation
$\mathit{eval}(g,V)\in\{\mathsf{true},\mathsf{false},\mathsf{unknown}\}$
is defined by the following absorbing rules: a comparison with a $\top$
operand evaluates to $\mathsf{unknown}$.  The connectives propagate
$\mathsf{unknown}$ as follows:
$\mathsf{unknown}\wedge\mathsf{false}=\mathsf{false}$,
$\mathsf{unknown}\wedge\mathsf{true}=\mathsf{unknown}$,
$\mathsf{unknown}\vee\mathsf{true}=\mathsf{true}$,
$\mathsf{unknown}\vee\mathsf{false}=\mathsf{unknown}$.
When a guard evaluates to $\mathsf{unknown}$ the transition is treated
as enabled.  For a branch, both successors are enabled simultaneously.
This over-approximation ensures no reachable state is missed at the
cost of exploring potentially infeasible paths.
\end{definition}

\begin{remark}[Place taxonomy and initial tokens]
\label{rem:places}
The three place classes and their initial token counts are:
\begin{enumerate}
  \item \textbf{Control places} $P_c$.  One place
    $\mathit{cp}(f,\mathit{sid})$ per statement per function, plus a
    return place $\mathit{cp}(f,\mathit{ret})$.  A token in
    $\mathit{cp}(f,\mathit{sid})$ means a thread is at statement
    $\mathit{sid}$ in $f$.  Initially, each spawned function's entry
    place $\mathit{cp}(f,\mathit{sid}^{0}_{f})$ holds one token.
    All other control places hold zero.

  \item \textbf{Resource places} $P_r$.  One place $\mathit{rp}(r)$
    per synchronization primitive.  Initial token counts:
    Mutex $\mapsto 1$,
    RwLock $\mapsto N$ where $N$ is the total number of concurrent
    entities,
    Semaphore$(n)\mapsto n$,
    Channel $\mapsto 0$,
    Condvar signal place $\mapsto 0$.
    \textsf{Var} and \textsf{Atomic} resources have no resource
    places.  Their state resides in the valuation $V$.

  \item \textbf{Auxiliary places} $P_a$.  For each \textsf{wait}
    site $\mathit{sid}$ on a condition variable: a \emph{wait place}
    $\mathit{wp}(\mathit{sid})$ (token = thread blocked) and a
    \emph{reacquire place} $\mathit{ra}(\mathit{sid})$ (token =
    thread woken, awaiting mutex).  All auxiliary places start at
    zero.
\end{enumerate}
\end{remark}

\begin{remark}[Condvar auxiliary variables]
\label{rem:condvar-vars}
Each condition variable $\mathit{cv}$ with wait sites
$\{w_1,\dots,w_k\}$ introduces auxiliary variables in $\mathit{Var}$:
a waiter-count variable $nw_{\mathit{cv}}\in\mathit{Var}$ with initial
value $0$, and a per-site notify-all flag
$na_{w_i}\in\mathit{Var}$ with initial value $\mathsf{false}$ for each
$i\in\{1,\dots,k\}$.
\end{remark}

\begin{remark}[Anchor mapping]
\label{rem:anchor}
$\mu$ is total on transitions generated from \textsc{Cir} statements.
For compound operations such as \textsf{wait}, which generates
multiple transitions (WaitEnter, Wake1, WakeA, Reacquire), each
transition is anchored to the originating \textsf{wait} statement's
$\mathit{sid}$.  Transitions synthesized from summaries are anchored
to $\mathit{sid}_{\bot}$ and excluded from witness traces.
\end{remark}

\begin{definition}[\textsc{Cvn} state]
\label{def:cvn-state}
A \textsc{Cvn} state is a pair $S=(M,V)$ where
$M:P\to\mathbb{N}_0$ is a marking and
$V:\mathit{Var}\to\mathit{Val}$ is a valuation.  The initial state
$S_0=(I_m,I_v)$ assigns token counts per Remark~\ref{rem:places} and
initial values per the \textsc{Cir} resource table $R$.  The pair
$(\mathcal{N},S_0)$ fully determines the reachable state space.
\end{definition}

\begin{definition}[Enabling and firing]
\label{def:enabling-firing}
A transition $t\in T$ is \emph{enabled} at state $S=(M,V)$ iff
\begin{enumerate}
  \item $\forall p\in P:\; M(p)\geq W_{\!\mathit{in}}(p,t)$, and
  \item $\mathit{eval}(G(t),V)\neq\mathsf{false}$.
\end{enumerate}
Upon firing, the successor state $S'=(M',V')$ is:
\[
  M'(p)=M(p)-W_{\!\mathit{in}}(p,t)+W_{\!\mathit{out}}(t,p)
  \quad\text{for each }p\in P,
\]
\[
  V'(x)=
  \begin{cases}
    \mathit{eval}\bigl(U(t)(x),\;V\bigr)
      & \text{if } x\in\mathrm{dom}(U(t)),\\
    V(x) & \text{otherwise.}
  \end{cases}
\]
All expressions in $U(t)$ are evaluated with respect to the
pre-firing valuation $V$.  If multiple transitions are simultaneously
enabled, the search explores all resulting interleavings.
\end{definition}

Each transition also carries a classification tag (Lock, Unlock,
Spawn, etc.) for diagnostics and visualization.  The tag does not
affect firing semantics.

\begin{definition}[Goal satisfaction]
\label{def:goal-sat}
A \textsc{Cvn} state $(M,V)$ \emph{satisfies} a \textsc{Cir} goal
$\gamma=(C_\gamma,V_\gamma)$, written $(M,V)\models\gamma$, iff every
completion and availability requirement holds in $M$ and
$\forall x\in\mathrm{dom}(V_\gamma):\;V(x)=V_\gamma(x)$.
A goal is \emph{reachable} if at least one reachable state satisfies
it.  The existential criterion is appropriate because concurrent
programs naturally admit multiple valid terminal configurations.  The
check answers: is the program still \emph{capable} of achieving its
intended outcome?
\end{definition}

% ═══════════════════════════════════════════════════════════════
\section{Translation Rules}
\label{app:translation-rules}

The translator $\mathit{translate}:\textsc{Cir}_v\to\textsc{Cvn}$ is
executed in three phases as described in
Section~\ref{subsec:translation}.  The detailed rules for Phase~2
(function-body translation) are given below.

\paragraph{Phase 1: Resource scanning.}
For each synchronization primitive in $R$, generate a resource place
$\mathit{rp}(r)$ and set its initial token count per
Remark~\ref{rem:places}.  Compute the RwLock concurrency bound $N$ as
the total number of concurrent entities.  For each \textsf{Var} and
\textsf{Atomic} resource, initialize the corresponding entry in the
valuation $I_v$.  For each \textsf{Condvar}, generate the auxiliary
places and variables described in Remarks~\ref{rem:places}
and~\ref{rem:condvar-vars}.

\paragraph{Phase 2: Function-body translation.}
For each function in $F$ and each statement in its body, generate
control places, transitions, weights, guards, and updates according
to Table~\ref{tab:translation-rules}.

\paragraph{Phase 3: Summary translation.}
For each function referenced by a \textsf{call} that has a summary in
$\Sigma$ but no body in $F$, generate a single atomic transition.
For each resource $x$ in the $\mathit{writes}$ set, the update sets
$x:=\top$.

The translation also maps each \textsc{Cir} goal
$\gamma=(C_\gamma,V_\gamma)$ to a \textsc{Cvn} reachability query:
$(f,\textsf{completed})$ maps to
$M(\mathit{cp}(f,\mathit{ret}))\geq 1$,
$(r,\textsf{available})$ maps to
$M(\mathit{rp}(r))\geq I_m(\mathit{rp}(r)))$, and variable
requirements pass through unchanged.  The mapping is purely internal:
users specify and receive results in \textsc{Cir} vocabulary.

% ═══════════════════════════════════════════════════════════════
\section{Bug Detection Details}
\label{app:bug-detection-details}

This appendix provides the full formal definitions of the bug
detection properties summarized in Section~\ref{subsec:bug-detection}.

\subsection{Deadlock Sub-Classification}

Two refinements specialize a deadlock diagnosis when the evidence is
unambiguous.  Each refinement is applied strictly \emph{after} a
deadlock has already been detected on a concrete witness, so it is
sound by construction: the set of reclassified states is a subset of
the set of deadlocks.

\begin{itemize}[leftmargin=1.5em]
  \item \emph{Signal-loss refinement $(\kappa_{sl})$.}
    A deadlock witness is relabelled $\kappa_{sl}$ when at least one
    blocked control token resides on a wait-place
    $wp_{cv,f,\mathit{sid}}$, or when the witness's firing sequence
    contains a \textsc{NotifyLost} (resp.\ \textsc{NotifyAllLost})
    transition that fired before the waiter entered its wait-place.
    The \textsc{NotifyLost} guard $nw_{cv}=0$ is evaluated on concrete
    integers only ($nw_{cv}$ is updated by literal $\pm 1$ steps), so
    the firing of this transition is detected precisely.  The
    refinement label is emitted only in conjunction with a confirmed
    deadlock, because an isolated lost notification is not itself a
    soundness violation on a scheduler-agnostic model.

  \item \emph{Channel-block refinement $(\kappa_{cb})$.}
    A deadlock witness is relabelled $\kappa_{cb}$ when at least one
    blocked consumer requires a token from a channel resource place
    whose marking is $0$.  Soundness reduces to the underlying
    deadlock predicate $\kappa_{dl}$.
\end{itemize}

\begin{remark}[Scope of soundness]
\label{rem:soundness-scope}
The \textsc{Cvn} state-space search reports two primary properties:
\emph{deadlock} and \emph{behavioral dead transition}, both sound
inside the model.  The sub-classifications $\kappa_{sl}$ and
$\kappa_{cb}$ are refinements of $\kappa_{dl}$ and inherit its
soundness.  The system never silently misses a deadlock or a dead
transition in the modelled \textsc{Cir}, and it never fabricates a bug
that is not already a genuine deadlock or a genuinely unreachable
transition under \textsc{Cvn} semantics.
\end{remark}

\subsection{Behavioral Dead Transition}

A transition $t\in T$ of the \textsc{Cvn} is \emph{behaviorally dead}
$(\kappa_{dt})$ with respect to the initial state $S_0$ iff
$t\notin \mathsf{fired}(\mathcal{R})$, where
$\mathsf{fired}(\mathcal{R})$ is the set of transitions that appear on
any edge of the reachability graph produced by exhaustive exploration
from $S_0$.  Intuitively, $t$ corresponds to a \textsc{Cir} statement
that the verifier can prove will never execute on any feasible
interleaving: either its guard is unsatisfiable, its input place never
receives tokens, or its predecessors are themselves dead.  Because
$\mathsf{fired}(\mathcal{R})$ is computed from the same reachability
graph used for deadlock detection, reporting $\kappa_{dt}$ requires no
additional search budget.

We report behavioral, not structural, dead transitions: a transition
that is structurally live but behaviorally dead (e.g., a branch that
the \textsc{Cvn}'s three-valued guard can statically refute) is still
reported.  Soundness is immediate: if $t$ fires on a real \textsc{Cir}
trace, then by Theorem~\ref{thm:forward} $t$ also appears on a
\textsc{Cvn} firing sequence, hence on an edge of $\mathcal{R}$.
Contrapositively, $t\notin\mathsf{fired}(\mathcal{R})$ implies that no
\textsc{Cir} execution ever reaches the statement anchored to $t$.
The conservative $\mathsf{unknown}$-branching only enlarges
$\mathsf{fired}(\mathcal{R})$, so it can only turn a would-be dead
transition into a live one, never the reverse.

\subsection{Goal Reachability Predicate Shapes}

Each business goal $\gamma\in\Gamma$ is compiled into a \emph{goal
predicate} $\phi_{\gamma}$ over \textsc{Cvn} states.  Three predicate
shapes are supported, chosen so that every field of the CIR goal
format maps onto a sound, marking-level query:

\begin{itemize}[leftmargin=1.5em]
  \item $\mathsf{Reachable}(p,k)\;\triangleq\;M(p)\ge k$ for
    control, wait, or resource places.  Used for ``thread $t$ reaches
    return'' ($p=cp_{t,\mathrm{ret}}$, $k=1$) and for ``at least $k$
    permits accumulated.''

  \item $\mathsf{Empty}(p)\;\triangleq\;M(p)=0$ for resource places
    whose initial marking is $0$ (channel message counters and condvar
    signal places).  This is the semantically meaningful availability
    check: asking $M(rp(r))\ge 0$ on a channel is trivially true and
    therefore worthless, whereas $\mathsf{Empty}(rp(r))$ expresses
    ``no pending messages.''

  \item $\mathsf{GlobalEq}(v,c)\;\triangleq\; V(v)=c$ for a variable
    $v$ in the \textsc{Cvn} variable store and a concrete value $c$.
    Because the \textsc{Cvn} uses a three-valued domain and
    $\mathsf{Unknown}$ never equals a concrete value, unresolved
    branches cannot spuriously satisfy a goal, preserving soundness.
\end{itemize}

A goal $\gamma$ with predicates
$\phi_{\gamma}^{1},\dots,\phi_{\gamma}^{n}$ is \emph{reachable},
written $\mathcal{R}\models\gamma$, iff some $s\in\mathcal{R}$
simultaneously satisfies all $\phi_{\gamma}^{i}$ (EF-reachability,
conjunctive).  The analyser returns the subset
$\mathcal{U}=\{\gamma\in\Gamma\mid\mathcal{R}\not\models\gamma\}$ of
unmet goals.

\section{Repair Diagnostics and Prompt Assembly}
\label{app:repair-details}

\subsection{Diagnostic Format}

\begin{definition}[Repair-oriented diagnostic]
\label{def:diagnostic}
A repair-oriented diagnostic is a tuple $\mathcal{D}=(\kappa,\;\pi_{\mu},\;\Sigma_{\mathit{state}},\;\Sigma_{\mathit{wait}},\;\Lambda,\;\Gamma_{\mathit{ctx}},\;H)$, where
$\kappa\in\{\kappa_{dl},\kappa_{sl},\kappa_{cb},\kappa_{dt}, \kappa_{goal}\}$ is the bug or goal kind, $\pi_{\mu}$ is a bug witness expressed as a sequence of
$\mathit{sid}$s via the anchor mapping $\mu$ (for a bug-free but goal-unmet report $\pi_{\mu}$ is empty), $\Sigma_{\mathit{state}}$ is a marking-derived state summary, $\Sigma_{\mathit{wait}}$ is a wait/ownership summary, $\Lambda$ is the relevant \textsc{Cir} slice, $\Gamma_{\mathit{ctx}}$ is the set of preservation constraints, and $H$ is a bug- or goal-specific repair hint.
\end{definition}

Each field serves a distinct repair purpose:

\begin{itemize}[leftmargin=1.5em]
  \item For deadlock and its sub-classifications, $\pi_{\mu}$ is a finite firing sequence ending in the bug state, expressed in \textsc{Cir} $\mathit{sid}$s rather than \textsc{Cvn} transition names.

  \item For a goal-unmet report ($\kappa=\kappa_{goal}$), $\pi_{\mu}$ is empty and $\Sigma_{\mathit{state}}$ instead lists each unmet predicate together with the \textsc{Cvn} place or variable it refers to.

  \item $\Sigma_{\mathit{state}}$ reports which thread is at which $\mathit{sid}$, which resource counts are available, and whether the system is globally blocked.

  \item $\Sigma_{\mathit{wait}}$ extracts the relations most relevant to repair: which thread holds which resource, which thread is waiting for which resource or condition, and which predicate triggered the classification.

  \item $\Lambda$ is a bounded window around the implicated $\mathit{sid}$s, restricting the LLM's editing scope.

  \item $\Gamma_{\mathit{ctx}}$ specifies what must survive the repair: resource names, thread structure, and business goals.

  \item $H$ offers a bug-specific suggestion such as enforce a consistent lock order or update the predicate before notification. For goal-unmet reports, it suggests the likely producer transition that must be re-introduced.
\end{itemize}

Together, the witness and state summary supply both \emph{causality} and \emph{locality}, which the raw marking alone cannot convey in a form the LLM can act on.

\subsection{Prompt Template}

\begin{table}[h]
\caption{Prompt template for \textsc{Cir} repair from a \textsc{Cvn}
diagnostic.}
\label{tab:repair-templates}
\centering\small
\begin{tabular}{@{}p{3.0cm}p{4.2cm}p{7.7cm}@{}}
\toprule
\textbf{Prompt field}
  & \textbf{Source}
  & \textbf{Repair role} \\
\midrule
Bug or goal kind
  & $\kappa$
  & Selects the repair mode: deadlock, signal-loss refinement, channel-block refinement, dead-transition, and goal-unmet each require different structural edits. \\
\addlinespace[3pt]
Witness trace
  & $\pi_{\mu}$
  & Localises the causal path to the bug state in editable $\mathit{sid}$s.  Empty for a goal-unmet report. \\
\addlinespace[3pt]
Bug-state summary
  & $\Sigma_{\mathit{state}}$
  & Reports where execution is stuck: current $\mathit{sid}$ per thread, blocked status,
    critical variable values. \\
\addlinespace[3pt]
Held resources
  & $\Sigma_{\mathit{wait}}$ (ownership)
  & Lists locks, channels, or semaphores currently preventing progress. \\
\addlinespace[3pt]
Waiting relations
  & $\Sigma_{\mathit{wait}}$ (waits)
  & Makes the dependency structure explicit: who waits for what, and which predicate triggered the diagnosis. \\
\addlinespace[3pt]
Relevant \textsc{Cir} slice
  & $\Lambda$
  & Restricts editing to a small region around implicated statements. \\
\addlinespace[3pt]
Preservation constraints
  & $\Gamma_{\mathit{ctx}}$
  & Prevents over-repair: resource names, thread structure, and business goals must survive. \\
\addlinespace[3pt]
Repair hint
  & $H$
  & Provides direction without forcing a single patch. \\
\addlinespace[3pt]
Output contract
  & Verifier schema
  & Forces return of a complete revised \textsc{Cir}, keeping the loop machine-checkable. \\
\bottomrule
\end{tabular}
\end{table}

\subsection{Goal-Violation Feedback}

\begin{figure}[h]
\begin{lstlisting}[basicstyle=\small\ttfamily,frame=single]
=== Goal Violation (Round 3) ===
Status: No concurrency bugs detected.

UNREACHABLE GOAL: G1
  desc: "Both threads complete, lock released,
         ready flag set"
  Missing: V[ready] = true
  No reachable state satisfies V[ready] = true.
  Hint: Ensure that a write(ready, true)
        operation exists on a reachable path.
\end{lstlisting}
\caption{Goal-violation feedback after bug-free verification.}
\label{fig:goal-violation}
\end{figure}

The prompt is intentionally compact.  It provides the minimum evidence needed for a local repair: what failed, how the execution reached the failure, what the system looks like at the failure point, and what must be preserved.


\section{Verification and Repair Hierarchy}
\label{app:repair-hierarchy}

\begin{table}[h]
\caption{Verification and repair hierarchy.}
\label{tab:two-layer}
\centering\footnotesize
\begin{tabular}{@{}p{5.8cm}ll@{}}
\toprule
\textbf{Issue} & \textbf{Layer} & \textbf{Tier} \\
\midrule
\multicolumn{3}{@{}l}{\textit{Structural (Layer 1)}} \\
\quad Missing unlock/drop       & \textsc{Cir} & 1 (auto-fix) \\
\quad Duplicate $\mathsf{sid}$  & \textsc{Cir} & 1 (auto-fix) \\
\quad Undefined resource        & \textsc{Cir} & 2 (LLM) \\
\quad Spawn/join mismatch       & \textsc{Cir} & 2 (LLM) \\
\quad Protection inconsistency  & \textsc{Cir} & 2 (LLM) \\
\midrule
\multicolumn{3}{@{}l}{\textit{Concurrency bugs (Layer 2a, deadlock search)}} \\
\quad Deadlock                                     & \textsc{Cvn} & 3 (LLM) \\
\quad \;\;$\hookrightarrow$ Signal-loss refinement & \textsc{Cvn} & 3 (LLM) \\
\quad \;\;$\hookrightarrow$ Channel-block refinement & \textsc{Cvn} & 3 (LLM) \\
\quad Cross-primitive deadlock                     & \textsc{Cvn} & 3 (LLM) \\
\quad Behavioral dead transition                    & \textsc{Cvn} & 3 (LLM) \\
\midrule
\multicolumn{3}{@{}l}{\textit{Functional regressions (Layer 2b, goal reachability)}} \\
\quad Unreachable business goal & \textsc{Cvn} & 3 (LLM) \\
\bottomrule
\end{tabular}
\end{table}

\end{document}